\begin{center}\begin{tikzpicture}[font=\small,>=stealth,line join=round]
\draw[gray!40,rounded corners=3pt] (0,0) rectangle (14.2,8.8);\node at (7.1,8.200000000000001) {例6：绝对值比较后，还要与两项定义域取交集};
\node[anchor=west] at (.6,6.8) {代数解 $x<0$ 或 $x>2$：};\draw[->,gray] (1,6.3)--(13,6.3);\draw[blue,very thick] (1.4,6.3)--(5,6.3);\draw[blue,very thick] (9,6.3)--(12,6.3);\draw[blue,fill=white] (5,6.3) circle(2.7pt);\node[below=7pt] at (5,6.3) {$0$};\draw[blue,fill=white] (9,6.3) circle(2.7pt);\node[below=7pt] at (9,6.3) {$2$};\node[anchor=west] at (.6,5.0) {两个输入都可代入：$-\frac12\le x\le1$：};\draw[->,gray] (1,4.5)--(13,4.5);\draw[violet,very thick] (4,4.5)--(7,4.5);\draw[violet,fill=violet] (4,4.5) circle(2.7pt);\node[below=7pt] at (4,4.5) {$-\frac12$};\draw[violet,fill=violet] (7,4.5) circle(2.7pt);\node[below=7pt] at (7,4.5) {$1$};\node[anchor=west] at (.6,3.1999999999999997) {最终 $[-\frac12,0)$：};\draw[->,gray] (1,2.6999999999999997)--(13,2.6999999999999997);\draw[orange!85!black,very thick] (4,2.6999999999999997)--(5,2.6999999999999997);\draw[orange!85!black,fill=orange!85!black] (4,2.6999999999999997) circle(2.7pt);\node[below=7pt] at (4,2.6999999999999997) {$-\frac12$};\draw[orange!85!black,fill=white] (5,2.6999999999999997) circle(2.7pt);\node[below=7pt] at (5,2.6999999999999997) {$0$};
\node[text width=12.6cm,align=center] at (7.1,.65) {三行共用同一 $x$ 尺度；右侧 $x>2$ 整段被定义域限制排除。};
\end{tikzpicture}\end{center}
